% =====================================================================
%  2bat-ud2-8.tex  —  SESSIÓ 8: Els vèrtexs i la funció objectiu
% =====================================================================
\horatitol{Sessió 8 — Els vèrtexs i la funció objectiu}%
{El cor de la programació lineal: mesurar allò que volem optimitzar i descobrir que l'òptim és sempre en un vèrtex.}

\subsection{Idea clau}

Ja sabem quines decisions són \emph{possibles} (la regió factible). Ara ve la pregunta
clau: entre totes, quina és la \textbf{millor}? Per respondre-la necessitem una
\textbf{funció objectiu} i un resultat sorprenent que simplifica molt la cerca.

\begin{idea}
\textbf{Funció objectiu.} És l'expressió que mesura allò que volem
\textbf{maximitzar} (cobertura, persones ateses, beneficis) o \textbf{minimitzar}
(cost), en funció de les dues variables. Per exemple:
\[
P=3x+2y \quad(\text{persones ateses}),\qquad C=4x+5y \quad(\text{cost en \euro}).
\]
\textbf{Resultat clau (teorema dels vèrtexs).} El valor òptim (màxim o mínim) d'una
funció objectiu sobre una regió factible s'assoleix \textbf{sempre en un vèrtex} de
la regió. Per tant:
\begin{center}
\textit{n'hi ha prou d'avaluar la funció objectiu a cada vèrtex i comparar.}
\end{center}
\textbf{Per què?} Els vèrtexs són les decisions «extremes», on s'esgota algun recurs.
La funció creix sempre cap a una direcció, així que el seu millor valor dins d'un
polígon cau en una cantonada, no al mig.
\end{idea}

\subsection{Exemples}

\begin{exemple}[avaluar una funció objectiu als vèrtexs]
Sobre la regió factible coneguda, de vèrtexs $(0,0)$, $(12,0)$, $(6,3)$ i $(0,5)$,
volem \textbf{maximitzar} les persones ateses $P=3x+2y$. Avaluem $P$ a cada vèrtex:
\begin{center}
\begin{tabular}{lcl}
\toprule
\textbf{Vèrtex} & $P=3x+2y$ & \\
\midrule
$(0,0)$ & $3\cdot 0+2\cdot 0=0$ & \\
$(12,0)$ & $3\cdot 12+2\cdot 0=36$ & $\leftarrow$ màxim \\
$(6,3)$ & $3\cdot 6+2\cdot 3=24$ & \\
$(0,5)$ & $3\cdot 0+2\cdot 5=10$ & \\
\bottomrule
\end{tabular}
\end{center}
El màxim és $P=36$, al vèrtex $(12,0)$: la decisió que atén més persones és fer
$12$ tallers i cap sortida.
\end{exemple}

\begin{exemple}[el mateix, però minimitzant]
Sobre la mateixa regió, si volguéssim \textbf{minimitzar} un cost $C=x+4y$,
avaluaríem $C$ als vèrtexs: $C(0,0)=0$, $C(12,0)=12$, $C(6,3)=18$, $C(0,5)=20$. El
mínim (a part del trivial $(0,0)$, que vol dir no fer res) seria al vèrtex de menys
cost útil. La idea és la mateixa: \textbf{avaluar als vèrtexs i comparar}, escollint el
més petit si minimitzem o el més gran si maximitzem.
\end{exemple}

\subsection{Exercicis resolts}

\begin{exresolt}[maximitzar avaluant als vèrtexs]
La regió factible té vèrtexs $(0,0)$, $(8,0)$, $(6,2)$ i $(0,5)$. Maximitza
$P=2x+3y$.
\solucio
Avaluem $P$ a cada vèrtex:
\[
P(0,0)=0,\quad P(8,0)=16,\quad P(6,2)=12+6=18,\quad P(0,5)=15.
\]
El més gran és $P=18$, al vèrtex $(6,2)$.
\resposta{Màxim $P=18$ al vèrtex $(6,2)$.}
\end{exresolt}

\begin{exresolt}[minimitzar un cost]
Sobre la regió de vèrtexs $(2,0)$, $(6,0)$, $(0,6)$ i $(0,3)$, minimitza el cost
$C=4x+5y$.
\solucio
Avaluem $C$ a cada vèrtex:
\[
C(2,0)=8,\quad C(6,0)=24,\quad C(0,6)=30,\quad C(0,3)=15.
\]
El més petit és $C=8$, al vèrtex $(2,0)$.
\resposta{Mínim $C=8$ al vèrtex $(2,0)$.}
\end{exresolt}

\subsection{Exercicis per fer}

\begin{exercicis}
\begin{exlist}
  \item Sobre la regió de vèrtexs $(0,0)$, $(10,0)$, $(4,6)$ i $(0,8)$, avalua
        $P=x+y$ a cada vèrtex i digues on és el màxim.
  \item Sobre la mateixa regió, avalua $P=3x+y$ a cada vèrtex i digues on és el màxim.
  \item Sobre la regió de vèrtexs $(0,0)$, $(6,0)$, $(6,3)$ i $(0,5)$, maximitza
        $P=2x+4y$.
  \item Sobre la regió de vèrtexs $(1,1)$, $(5,1)$, $(5,4)$ i $(1,4)$, minimitza el
        cost $C=2x+3y$.
  \item Una funció objectiu val $P=5x+5y$. Sobre la regió de vèrtexs $(0,0)$, $(4,0)$,
        $(2,2)$ i $(0,4)$, avalua-la a cada vèrtex. Què observes en dos d'ells?
  \item \textbf{(Raonament)} Explica amb les teves paraules per què, per trobar el
        màxim d'una funció objectiu, no cal provar tots els (infinits) punts de la
        regió factible, sinó només els vèrtexs.
\end{exlist}
\end{exercicis}

% ---------------------------------------------------------------------
\fullsolucions{Sessió 8}
\begin{solucions}
\begin{sollist}
  \item $P(0,0)=0$, $P(10,0)=10$, $P(4,6)=10$, $P(0,8)=8$. Màxim $P=10$, assolit a
        $(10,0)$ i a $(4,6)$ (empat).
  \item $P(0,0)=0$, $P(10,0)=30$, $P(4,6)=18$, $P(0,8)=8$. Màxim $P=30$ a $(10,0)$.
  \item $P(0,0)=0$, $P(6,0)=12$, $P(6,3)=24$, $P(0,5)=20$. Màxim $P=24$ a $(6,3)$.
  \item $C(1,1)=5$, $C(5,1)=13$, $C(5,4)=22$, $C(1,4)=14$. Mínim $C=5$ a $(1,1)$.
  \item $P(0,0)=0$, $P(4,0)=20$, $P(2,2)=20$, $P(0,4)=20$. Tres vèrtexs donen $20$:
        la funció és constant sobre la vora que els uneix (hi ha múltiples òptims).
  \item Resposta oberta. Idea: la funció objectiu creix de manera uniforme en una
        direcció; sobre un polígon, el seu valor extrem s'assoleix sempre en una
        cantonada (vèrtex). Per això n'hi ha prou d'avaluar els vèrtexs, que són un
        nombre finit, en lloc dels infinits punts interiors.
\end{sollist}
\end{solucions}
